%\ifodd\value{page}\hbox{}\newpage\fi
\chapter*{Dhyanam 3}

\begin{sloka}
{\sanskritfont
\begin{tabular}{c}
ध्यायेत् पद्मासनस्थां विकसितवदनां पद्मपत्रायताक्षीं \\
हेमाभां पीतवस्त्रां करकलितलसद्धेमपद्मां वराङ्गीम् ।\\
सर्वालङ्कारयुक्तां सततमभयदां भक्तनम्रां भवानीं \\
श्रीविद्यां शान्तमूर्तिं सकलसुरनुतां सर्वसम्पत्प्रदात्रीम् ॥
\end{tabular}

\vspace{1.5em}

{\small \iastfont
dhyāyet padmāsanasthāṃ vikasitavadanāṃ padmapatrāyatākṣīm \\
hemābhāṃ pītavastrāṃ karakalitalasaddhemapadmāṃ varāṅgīm |\\
sarvālaṅkārayuktāṃ satatamabhayadāṃ bhaktanamrāṃ bhavānīṃ \\
śrīvidyāṃ śāntamūrtiṃ sakalasuranutāṃ sarvasampatpradātrīm ||\\
}

\vspace{1em}

{\small
\anvaya{%
पद्मासनस्थाम् विकसितवदनाम् पद्मपत्रायताक्षीम्
हेमाभाम् पीतवस्त्राम् करकलितलसद्धेमपद्माम् वराङ्गीम् ।
सर्वालङ्कारयुक्ताम् सततम् अभयदाम् भक्तनम्राम् भवानीम्
श्रीविद्याम् शान्तमूर्तिम् सकलसुरनुताम् सर्वसम्पत्प्रदात्रीम् ध्यायेत् ॥
}
}

\small \iastfont \itmean{%
«Si mediti su Bhavānī, seduta sul loto, dal volto aperto e luminoso,
dagli occhi simili a petali di loto, dal fulgore dorato e vestita di giallo,
che tiene in mano uno splendente loto d’oro, ornata di ogni gioiello,
benevola verso i devoti, sempre dispensatrice di protezione,
incarnazione della Śrīvidyā, forma di pace, lodata da tutti gli dèi,
dispensatrice di ogni pienezza.»%
}


}
\end{sloka}

\wordmeaning{%
\wordline{ध्यायेत्}{dhyāyet}
  { si mediti;}%
\wordline{पद्मासनस्थाम्}{padmāsanasthām}
  { seduta sul loto;}%
\wordline{विकसितवदनाम्}{vikasitavadanām}
  { dal volto luminoso e aperto;}%
\wordline{पद्मपत्रायताक्षीम्}{padmapatrāyatākṣīm}
  { dagli occhi simili a petali di loto;}%
\wordline{हेमाभाम्}{hemābhām}
  { dal fulgore dorato;}%
\wordline{पीतवस्त्राम्}{pītavastrām}
  { vestita di giallo;}%
\wordline{करकलितलसद्धेमपद्माम्}{karakalitalasaddhemapadmām}
  { che tiene in mano uno splendente loto d’oro;}%
\wordline{वराङ्गीम्}{varāṅgīm}
  { dalla forma eccelsa;}%
\wordline{सर्वालङ्कारयुक्ताम्}{sarvālaṅkārayuktām}
  { ornata di ogni gioiello;}%
\wordline{सततम् अभयदाम्}{satatam abhayadām}
  { che concede costantemente protezione;}%
\wordline{भक्तनम्राम्}{bhaktanamrām}
  { benevola verso i devoti umili;}%
\wordline{श्रीविद्याम्}{śrīvidyām}
  { incarnazione della Śrīvidyā;}%
\wordline{शान्तमूर्तिम्}{śāntamūrtim}
  { dalla forma pacificata;}%
\wordline{सकलसुरनुताम्}{sakalasuranutām}
  { lodata da tutti gli dèi;}%
\wordline{सर्वसम्पत्प्रदात्रीम्}{sarvasampatpradātrīm}
  { dispensatrice di ogni pienezza;}%
}

